% Part: first-order-logic
% Chapter: models-theories
% Section: expressing-props-of-structures

\documentclass[../../../include/open-logic-section]{subfiles}

\begin{document}

\olfileid[fr]{fol}{mat}{exs}

\olsection{Exprimer des propriétés des structures}

\begin{explain}
Il est souvent utile et important d'énoncer des conditions portant
sur des fonctions et des relations, ou, plus généralement, de dire
que les fonctions et les relations d'une structure satisfont à ces
conditions. Nous voudrions, par exemple, distinguer les
!!{structure}s d'un langage qui « rendent compte » du sens que nous
voulons donner aux !!{predicate}s de celles qui ne le font pas.
Nous sommes bien sûr entièrement libres de préciser quelles
!!{structure}s nous « visons » : nous pouvons exiger que
l'interprétation du !!{predicate}~$\le$ soit un ordre, ou ne nous
intéresser qu'aux interprétations de~$\Lang L$ dont le domaine est
formé d'ensembles et où $\Obj \in$ désigne la relation
« est !!a{element} de ». Mais pouvons-nous formuler ces exigences
au moyen d'!!{sentence}s du langage ? En d'autres termes, quelles
conditions sur !!a{structure}~$\Struct M$ peut-on exprimer par
!!a{sentence}, ou éventuellement par un ensemble d'!!{sentence}s,
du langage de~$\Struct M$ ? Certaines conditions ne peuvent pas
l'être. Par exemple, aucun énoncé satisfaisable de~$\Lang L_A$ ne peut imposer\footnote{Précision éditoriale : nous excluons les énoncés insatisfaisables, pour lesquels toute restriction sur les modèles serait vraie par vacuité.}
qu'une !!{structure}~$\Struct M$ dans laquelle il est vrai ait
$\Domain M = \Nat$. Nous ne pouvons pas exprimer
« le domaine ne contient que les nombres naturels ». Mais nous
pouvons peut-être exprimer certaines « propriétés structurelles »
des !!{structure}s. Lesquelles peut-on exprimer par des
!!{sentence}s ? Ou, ce qui revient au même, quelles collections de
!!{structure}s peut-on décrire comme celles qui satisfont
!!a{sentence}, ou un ensemble d'!!{sentence}s ?
\end{explain}

\begin{defn}[Modèle d'un ensemble]
Soit $\Gamma$ un ensemble d'!!{sentence}s d'un langage~$\Lang L$.
Nous disons qu'!!a{structure}~$\Struct M$ \emph{est un modèle
de}~$\Gamma$ si $\Sat{M}{!A}$ pour tout $!A \in \Gamma$.
\end{defn}

\begin{ex}
L'énoncé $\lforall[x][x \le x]$ est vrai dans~$\Struct M$ si et
seulement si $\Assign{\le}{M}$ est une relation réflexive. L'énoncé
$\lforall[x][\lforall[y][((x \le y \land y \le x) \lif x = y)]]$
est vrai dans~$\Struct M$ si et seulement si $\Assign{\le}{M}$
est antisymétrique. L'énoncé
$\lforall[x][\lforall[y][\lforall[z][((x \le y \land y \le z)
      \lif x \le z)]]]$ est vrai dans~$\Struct M$ si et seulement si
$\Assign{\le}{M}$ est transitive. Les modèles de
\begin{align*}
\{\quad &\lforall[x][x \le x], \\
   & \lforall[x][\lforall[y][((x \le y \land y \le
    x) \lif x = y)]], \\
   &\lforall[x][\lforall[y][\lforall[z][((x \le y
      \land y \le z) \lif x \le z)]]] \quad \}
\end{align*}
sont donc exactement les structures dans lesquelles
$\Assign{\le}{M}$ est réflexive, antisymétrique et transitive,
c'est-à-dire un ordre partiel. Nous pouvons donc prendre ces énoncés
comme axiomes de la \emph{théorie du premier ordre des ordres
partiels}.
\end{ex}

\end{document}
